
\documentclass[11pt]{article}
\usepackage[english]{babel}
\usepackage[a4paper,margin=0.9in]{geometry}
\usepackage[T1]{fontenc}
\usepackage[utf8]{inputenc}
\usepackage{lmodern}
\usepackage{amsmath,amssymb,amsthm,mathtools,mathrsfs}
\usepackage{microtype}
\usepackage{fancyhdr}
\usepackage{array,longtable,enumitem,multicol}
\usepackage{tikz}
\setlength{\parindent}{1.5em}
\setlength{\parskip}{0.18em}
\setlength{\headheight}{14pt}
\emergencystretch=8em
\pagestyle{fancy}
\fancyhf{}
\cfoot{\thepage}
\lhead{Heinrich Weber, Lehrbuch der Algebra, Vol. II}
\newcommand{\eps}{\varepsilon}
\newcommand{\ii}{\mathrm{i}}
\newcommand{\Om}{\Omega}
\newcommand{\tg}{\operatorname{tg}}
\newcommand{\PSL}{\operatorname{PSL}}
\newcommand{\Jac}{\operatorname{Jac}}
\newcommand{\mat}[1]{\begin{pmatrix}#1\end{pmatrix}}
\newcommand{\arr}[1]{\begin{array}{#1}}
\newcommand{\G}{G_{168}}
\newcommand{\Pgrp}{P_{168}}
\newcommand{\slfrac}[2]{{#1}/{#2}}
\rhead{English Translation}
\begin{document}

\section*{Fifteenth Section. The form problem of the group \(G_{168}\) and the theory of equations of the seventh degree.}

\section*{\S\,125. The resolvents of the form problem.}

For the group \(G_{168}\), the form problem gives, according to the general theory, an equation of degree 168 whose coefficients are invariants. Each divisor of the group gives a resolvent of lower degree, namely of degree equal to the index of the divisor. We shall use here the two most interesting cases, the octahedral group
\begin{equation}
G_{24}=\langle \chi^2,\omega,\Theta^\nu\rangle,
\tag{1}
\end{equation}
and the group of degree 21
\begin{equation}
G_{21}=\langle \tau,\Theta\chi\rangle,
\tag{2}
\end{equation}
which lead to resolvents of the seventh and eighth degrees.

For the resolvents of the seventh degree one may start from the principal axes of the group. A principal axis remains fixed under the group \(G_{24}\), and under the whole group is carried into three distinct lines. Thus there are seven triples of principal axes, determined by an equation of the seventh degree. With the notation of \S\,118 put
\begin{equation}
\begin{aligned}
A_1&=\alpha\gamma x_1+\gamma\beta\eps^3x_2+\alpha\beta\eps^2x_3,\\
A_2&=\beta\gamma\eps^3x_1+\alpha\beta\eps^2x_2+\alpha\gamma x_3,\\
A_3&=\alpha\beta\eps^2x_1+\alpha\gamma x_2+\beta\gamma\eps^3x_3.
\end{aligned}
\tag{3}
\end{equation}
Then \(A_1=0,A_2=0,A_3=0\) are the equations of three of these axes. The substitution \(\chi\) carries \(A_1,A_2,A_3\) cyclically into one another, and the other generating substitutions merely permute them. Therefore the symmetric functions of these three quantities are roots of resolvents of degree seven.

The simplest such symmetric function is
\begin{equation}
A_1^2+A_2^2+A_3^2.
\tag{4}
\end{equation}
Using the relations among \(\alpha,\beta,\gamma\) one finds
\begin{equation}
\mu={2\eps\over7},\qquad
\lambda=-\eps {1+\sqrt{-7}\over14},\qquad {\mu\over\lambda}=-{1-\sqrt{-7}\over2}.
\tag{5}
\end{equation}
Thus, if \(A_1^2+A_2^2+A_3^2=\lambda z\),
\begin{equation}
z=x_1^2+x_2^2+x_3^2-{1-\sqrt{-7}\over2}(x_2x_3+x_3x_1+x_1x_2).
\tag{6}
\end{equation}
The other roots are obtained by the powers of \(\tau\):
\begin{equation}
\begin{aligned}
z_r={}&\eps^{2r}x_1^2+\eps^{4r}x_2^2+\eps^rx_3^2\\
&-{1-\sqrt{-7}\over2}\bigl(\eps^{-r}x_2x_3+\eps^{-2r}x_3x_1+\eps^{-4r}x_1x_2\bigr),\quad r=0,1,\ldots,6.
\end{aligned}
\tag{7}
\end{equation}
The coefficients of the equation whose roots are the seven quantities \(z_r\) are invariants. If \(a_\nu\) is the coefficient of \(z^{7-\nu}\), then \(a_\nu\) is an invariant of degree \(2\nu\). Hence \(a_1=0\). The other coefficients are linear and homogeneous combinations of
\begin{equation}
\begin{array}{c|c}
a_2&f\\
a_3&\Delta\\
a_4&f^2\\
a_5&\Delta f\\
a_6&\Delta^2,\ f^3\\
a_7&C,\ \Delta f^2.
\end{array}
\tag{8}
\end{equation}
In the special case \(f=0\), comparison of leading terms in the power sums gives
\begin{equation}
\sum z_r^3=-3\cdot7\,{1-\sqrt{-7}\over2}x_1^5x_2+\cdots,
\end{equation}
\begin{equation}
\sum z_r^6=\left[6\cdot7+15\cdot7\left({1-\sqrt{-7}\over2}\right)^2\right]x_1^{10}x_2^2+\cdots.
\end{equation}
Newton's formulae then give
\begin{equation}
z^7-7{1-\sqrt{-7}\over2}\Delta z^4-7{5+\sqrt{-7}\over2}\Delta^2z-C=0.
\tag{9}
\end{equation}
With
\begin{equation}
z={Cu\over\Delta^2},\qquad g={\Delta^7\over C^3},
\tag{10}
\end{equation}
this becomes
\begin{equation}
u^7-7{1-\sqrt{-7}\over2}g u^4-7{5+\sqrt{-7}\over2}g^2u-g^2=0.
\tag{11}
\end{equation}
In the general case another parameter is required,
\begin{equation}
h={f\Delta^4\over C^2}.
\tag{12}
\end{equation}
The coefficients of the general resolvent are, apart from numerical factors,
\begin{equation}
h,\ g,\ h^2,\ hg,\ (g^2,h^3),\ (g^2,h^2g),
\end{equation}
and direct comparison gives
\begin{equation}
\begin{array}{rcl}
&& 7{1-\sqrt{-7}\over2}h,\\
&&-7{1-\sqrt{-7}\over2}g,\\
&&-7(4+\sqrt{-7})h^2,\\
&&14(2+\sqrt{-7})hg,\\
&&-7{5+\sqrt{-7}\over2}g^2-7{7+3\sqrt{-7}\over2}h^3,\\
&&-g^2+{167-7\sqrt{-7}\over2}gh^2.
\end{array}
\tag{13}
\end{equation}
The functions \(g,h\) are fractional invariants which depend only on the ratios \(x_1:x_2:x_3\). Every fractional invariant with the same property is rationally expressible through \(g,h\). Indeed, its numerator and denominator are whole invariants of the same degree, and therefore are rationally expressible in \(f,\Delta,C\). If a term has the form
\begin{equation}
f^a\Delta^b C^c,
\end{equation}
then, for common degree \(m\),
\begin{equation}
4a+6b+14c=m.
\tag{14}
\end{equation}

\section*{\S\,126. Reduction of the general resolvent of the seventh degree to the special one.}

For a point \((x)\) write
\begin{equation}
u_x={\Delta_x^2\over C_x},\qquad g_x={\Delta_x^7\over C_x^3},\qquad h_x={f_x\Delta_x^4\over C_x^2}.
\tag{1}
\end{equation}
The general resolvent is
\begin{equation}
R(u_x,g_x,h_x)=R_x=0,
\tag{2}
\end{equation}
and the special resolvent is obtained by setting \(f_x=0\).

According to Klein, the solution of the general resolvent can be reduced to that of the special resolvent if one adjoins a root of a biquadratic equation. Introduce a second point \((y)\), and form the polar of \(f\) with respect to \((x)\):
\begin{equation}
f_1(x,y)=y_1f_{x_1}(x)+y_2f_{x_2}(x)+y_3f_{x_3}(x).
\tag{3}
\end{equation}
Require \((y)\) to lie both on the ground curve and on this polar,
\begin{equation}
f(y)=0,
\qquad
f_1(x,y)=0.
\tag{4}
\end{equation}
Then to a point \((x)\) there correspond four points \((y)\). When \((x)\) is replaced by a connected point under the group, these four points also go into connected points. For such a point \((y)\), one has \(h_y=0\); the quantity \(g_y\), however, is a four-valued function of \((x)\). Its symmetric functions are invariants of \(G_{168}\), and hence \(g_y\) satisfies a biquadratic equation with coefficients rational in \(g_x,h_x\).

Substituting the corresponding value \(u_y\) into the special resolvent gives
\begin{equation}
\Phi(u_y,u_x,g_x,h_x)=0.
\tag{5}
\end{equation}
If \((x)\) is chosen so that no two of the four points \(y',y'',y''',y''''\) are connected, then \(\Phi=0\), considered as an equation in \(u_x\), has only one root in common with \(R_x=0\). The greatest common divisor of \(R\) and \(\Phi\) therefore gives \(u_x\) rationally in terms of \(u_y,g_x,h_x\). This is the desired reduction.

The same group contains a divisor of degree 21, generated by \(\tau\) and \(\chi\), which gives a resolvent of degree eight. One may take \(x_1x_2x_3\) as a root. Its coefficients are invariants rising to degree 94.

\section*{\S\,127. The permutation group of degree 168.}

The group of the equations obtained above is the group of ternary linear congruence substitutions modulo 2. The seven non-zero triples are
\begin{equation}
(1,0,0),(0,1,0),(0,0,1),(0,1,1),(1,0,1),(1,1,0),(1,1,1).
\tag{1}
\end{equation}
A substitution
\begin{equation}
\begin{pmatrix}
a&b&c\\ a_1&b_1&c_1\\ a_2&b_2&c_2
\end{pmatrix}
\tag{2}
\end{equation}
with entries taken modulo 2 permutes these seven triples, and the totality of these permutations is \(P_{168}\). With
\begin{equation}
\tau=\begin{pmatrix}1&0&1\\1&0&0\\0&1&0\end{pmatrix}
\tag{3}
\end{equation}
and with \(z_0=(1,0,0)\), the successive powers of \(\tau\) arrange the seven symbols cyclically. The two generators are even permutations, and therefore \(P_{168}\) is a subgroup of the alternating group on seven letters.

\section*{\S\,128. Equations of the seventh degree with a group of degree 168.}

Let \(\lambda_0,\lambda_1,\ldots,\lambda_6\) be seven quantities and let
\begin{equation}
\tau=(\lambda_0,\lambda_1,\lambda_2,\lambda_3,\lambda_4,\lambda_5,\lambda_6)
\tag{1}
\end{equation}
be a cycle of degree seven. Add
\begin{equation}
\omega=(\lambda_1,\lambda_2)(\lambda_3,\lambda_5).
\tag{2}
\end{equation}
These permutations generate \(P_{168}\). In this notation
\begin{equation}
\Theta^3=\tau^6\omega\tau\omega=(1,2,5,3)(4,6),\qquad
\Theta=(1,3,5,2)(4,6),
\end{equation}
\begin{equation}
\chi=\omega\tau^3\Theta\tau^2=(1,4,2)(3,5,6).
\end{equation}
The element \(\lambda_0\) is fixed by the octahedral subgroup \(P_{24}\). A function belonging to \(P_{168}\) is obtained by taking
\begin{equation}
\begin{aligned}
v={}&\lambda_0\lambda_2\lambda_6+
\lambda_1\lambda_5\lambda_0+
\lambda_2\lambda_6\lambda_1+
\lambda_3\lambda_0\lambda_2\\
&+\lambda_4\lambda_1\lambda_3+
\lambda_5\lambda_2\lambda_4+
\lambda_6\lambda_3\lambda_5.
\end{aligned}
\tag{3}
\end{equation}
This remains unchanged under \(\omega\) and \(\tau\), hence under all of \(P_{168}\). Consequently \(v\) is a root of a degree-thirty equation with symmetric coefficients; after adjoining the square-root of the discriminant it splits into two factors of degree fifteen.

If \(v\) belongs to the rationality domain, together with the symmetric functions of the \(\lambda\), then the \(\lambda\)'s are roots of a special equation of the seventh degree whose Galois group has order 168. It remains to form three rational functions \(X_1,X_2,X_3\) of the roots which transform according to the corresponding ternary substitutions of \(G_{168}\). Substituting these functions in the invariants of \(G_{168}\) then gives quantities in the rationality domain, and any function of \(X_1,X_2,X_3\) changed by \(G_{168}\) is a root of a total resolvent.

\section*{\S\,129. Contragredient groups.}

If
\begin{equation}
y=A(x),\qquad \eta=A^{-t}(\xi),
\tag{1}
\end{equation}
are transposed substitutions, then the rows \((x_1,x_2,x_3)\) and \((\xi_1,\xi_2,\xi_3)\) are contragredient. For derivatives this gives the rule
\begin{equation}
\xi_1^\kappa\xi_2^\lambda\xi_3^\mu
=\sum C^{\kappa\lambda\mu}_{\alpha\beta\gamma}
\eta_1^\alpha\eta_2^\beta\eta_3^\gamma,
\qquad \kappa+\lambda+\mu=m,
\tag{2}
\end{equation}
and hence
\begin{equation}
{\partial^m \varphi\over\partial x_1^\kappa\partial x_2^\lambda\partial x_3^\mu}
=\sum C^{\kappa\lambda\mu}_{\alpha\beta\gamma}
{\partial^m \varphi\over\partial y_1^\alpha\partial y_2^\beta\partial y_3^\gamma}.
\tag{3}
\end{equation}
This proves the transformation law for higher derivatives. Therefore powers and products of variables in a contravariant may be replaced by the corresponding differential operators applied to a covariant, and conversely. If \(A\) runs through a group \(G\), then \(A^{-t}\) runs through the contragredient group. The invariants of the latter are the contravariants of \(G\).

\section*{\S\,130. Solution of the seventh-degree equation with group \(P_{168}\) by the form problem of \(G_{168}\).}

The group \(G_{168}\) is contragredient to itself. Its generators remain in the group under transposition. Hence its contravariants coincide with its invariants, up to notation. The octahedral subgroup \(G_{24}\), however, is not equal to its contragredient.

Let \(\eta_0\) be a function of the roots \(\lambda_0,\ldots,\lambda_6\) belonging to \(P_{24}\), for instance a rational function of \(\lambda_0\), and let \(\eta_0,\eta_1,\ldots,\eta_6\) be its cyclic conjugates. With the forms
\begin{equation}
z_r=\eps^{2r}x_1^2+\eps^{4r}x_2^2+\eps^rx_3^2-{1-\sqrt{-7}\over2}(\eps^{-r}x_2x_3+\eps^{-2r}x_3x_1+\eps^{-4r}x_1x_2)
\tag{1}
\end{equation}
put
\begin{equation}
\psi=\sum_{r=0}^6\eta_r z_r.
\tag{2}
\end{equation}
Then
\begin{equation}
\psi=p_1x_1^2+p_2x_2^2+p_3x_3^2+2q_1x_2x_3+2q_2x_3x_1+2q_3x_1x_2.
\tag{3}
\end{equation}
Taking three such functions gives three forms \(\psi,\psi',\psi''\). Their functional determinant
\begin{equation}
\Psi=
\begin{vmatrix}
\psi_{x_1}&\psi_{x_2}&\psi_{x_3}\\
\psi'_{x_1}&\psi'_{x_2}&\psi'_{x_3}\\
\psi''_{x_1}&\psi''_{x_2}&\psi''_{x_3}
\end{vmatrix}
\tag{4}
\end{equation}
is a cubic form
\begin{equation}
\Psi=\sum A_{ijk}x_ix_jx_k.
\tag{5}
\end{equation}
Its coefficients are trilinear functions of the three systems \(\eta_r,\eta'_r,\eta''_r\). By using contragredient variables and then applying the theorem on substituting derivatives in invariants, one obtains three functions
\begin{equation}
X_1,
\qquad X_2,
\qquad X_3,
\tag{6}
\end{equation}
which undergo exactly the linear substitutions corresponding to the permutations of \(P_{168}\). This is the required connection with the form problem.

\section*{\S\,131. Possibility of determining the functions \(X_1,X_2,X_3\).}

It remains to show that \(\eta_0,\eta'_0,\eta''_0\) can be chosen so that \(X_1,X_2,X_3\) are not identically zero. From the calculation of the preceding paragraph one obtains expressions of the following type:
\begin{equation}
\begin{aligned}
X_1&=(p_3,p'_2,q''_1)+(q_1,p'_1,q''_2)+(p_1,q'_3,p''_3),\\
X_2&=(q_2,q'_1,p''_3)+(p_3,p'_2,q''_2)+(q_3,p'_2,q''_1),\\
X_3&=(q_2,q'_3,q''_2)+(p_3,q'_1,p''_3)+(q_2,p'_2,p''_3),
\end{aligned}
\tag{1}
\end{equation}
where the bracket denotes the determinant formed from the indicated rows of coefficients. These functions are trilinear in the three chosen systems.

Make the substitution
\begin{equation}
\eta_r=\sum_s a_s\lambda_r^s,
\qquad
\eta'_r=\sum_s a'_s\lambda_r^s,
\qquad
\eta''_r=\sum_s a''_s\lambda_r^s.
\tag{2}
\end{equation}
The determinant of this transformation is, up to sign, the product of all differences \(\lambda_i-\lambda_k\), and is non-zero when the roots are distinct. Hence rational values of \(a_s,a'_s,a''_s\) can be chosen so that the three trilinear forms do not vanish. The required functions \(X_1,X_2,X_3\) therefore exist, and the reduction of the seventh-degree equation with group \(P_{168}\) to the form problem of \(G_{168}\) is complete.

\end{document}
